\documentclass{article}
\usepackage[UTF8]{ctex}
\usepackage{amsmath,amssymb,amsfonts}
\usepackage{enumitem}
\usepackage{geometry}
\geometry{a4paper, margin=1in}
% === 错题本排版优化（不参与题目标识） ===
\setlength{\parindent}{0pt}
\setlength{\parskip}{0.35em plus 0.10em minus 0.08em}
\setlength{\emergencystretch}{3em}
\setlist{itemsep=0.25em, topsep=0.40em, parsep=0pt}
\allowdisplaybreaks[2]
% === 错题本排版优化结束 ===
\begin{document}
\section*{国防科技大学2026年高等代数真题}
\section*{142．国防科技大学2026年数学分析与高等代数试卷（回忆版）}
（注．题目来自考生回忆，部分表述或数据可能有误）

\subsection{题目1}
1．令 \(D(x)=\left\{\begin{array}{l}1, x \text { 为有理数，} \\ 0, x \text { 为无理数。 }\end{array}\right.\)

\begin{itemize}
  \item[（1）] 证明：\(D(x)\) 在实数范围内处处不存在极限．
  \item[（2）] 构造 \(f(x)\) ，使其在 \(x=1\) 和 \(x=2\) 处连续，其余点处均不连续．
\end{itemize}

\subsection{题目2}
2．已知 \(A^{3}=2 E, B=A^{2}-2 A+2 E\) ，求 \(B^{-1}\) ．

\subsection{题目3}
3．给定 \(P^{3}\) 的两组基

\[
\begin{gathered}
\alpha_{1}=(1,0,1), \alpha_{2}=(2,1,0), \alpha_{3}=(1,1,1) \\
\beta_{1}=(1,2,-1), \beta_{2}=(2,2,-1), \beta_{3}=(2,-1,-1) .
\end{gathered}
\]

线性变换 \(\sigma\) 满足 \(\sigma\left(\alpha_{i}\right)=\beta_{i}(i=1,2,3)\) 。

\begin{itemize}
  \item[（1）] 求由基 \(\alpha_{1}, \alpha_{2}, \alpha_{3}\) 到 \(\beta_{1}, \beta_{2}, \beta_{3}\) 的过渡矩阵。
  \item[（2）] 写出 \(\sigma\) 在基 \(\alpha_{1}, \alpha_{2}, \alpha_{3}\) 下的矩阵。
\end{itemize}

\subsection{题目4}
4．未知。

\subsection{题目5}
5．计算 \(\iint_{D}\left|(x-1)^{2}+(y-1)^{2}-1\right| \mathrm{d} x \mathrm{~d} y\) ，其中 \(D=\left\{(x, y) \mid x^{2}+y^{2} \leq 9\right\}\) ．

\subsection{题目6}
6．令 \(A=\left(\begin{array}{ccc}1 & -3 & 1 \\ 2 & 1 & 0 \\ 3 & 1 & 1\end{array}\right)\) ．

\begin{itemize}
  \item[（1）] 证明：\(A\) 在有理数域上不可以对角化．
  \item[（2）] 证明：\(A\) 在复数域上可以对角化．
\end{itemize}

\subsection{题目7}
7．设 \(f(x)=x^{m}\)（ \(m\) 为大于 2 的正整数），证明：\(f(x)\) 在 \([0,+\infty)\) 上不一致连续．

\subsection{题目8}
8．设 \(f(x)\) 在 \([0,1]\) 上连续可导，且 \(f(0)=0\) ，证明： \(\int_{0}^{1}\left|f(x) f^{\prime}(x)\right| \mathrm{d} x \leq \frac{1}{2} \int_{0}^{1}\left[f^{\prime}(x)\right]^{2} \mathrm{~d} x\) ．

\subsection{题目9}
9．证明：\(a_{n}=\sum_{k=1}^{n} \frac{1}{n+k}\) 收敛。

\subsection{题目10}
10．设 \(n\) 阶实对称矩阵 \(A\) 的全部特征值按大小顺序排列成 \(\lambda_{1} \geq \lambda_{2} \geq \cdots \geq \lambda_{n}\) ，证明：

\[
\max _{x \neq 0} \frac{x^{\mathrm{T}} A x}{x^{\mathrm{T}} x}=\lambda_{1}, \min _{x \neq 0} \frac{x^{\mathrm{T}} A x}{x^{\mathrm{T}} x}=\lambda_{n} .
\]
\end{document}
